% GENERATED FILE. Edit canonical lessons, then run npm run generate:books.
% canonical-source-hash: fnv1a64:059f46566c6556b0
% canonical-lessons: ZH-C507-xiongdi, ZH-C507-fuzi, ZH-W507-fu-father, ZH-C507-funu, ZH-C507-fuqi, ZH-C507-zhangbei

\chapter{Fathers and Sons}
\label{ch:zh-fathers-and-sons}
\hlchaptermodality{507}

\begin{chapteropening}
\textbf{By the end of this chapter:} \emph{I can talk about brothers, fathers and sons, fathers and daughters, married couples and my elders.}

Complete the last lesson of chapter 507: I can talk about brothers, fathers and sons, fathers and daughters, married couples and my elders.
\end{chapteropening}

\section[xiōngdì]{xiōngdì — brothers}

\label{lesson:ZH-C507-xiongdi}



\begin{quote}
Before the new one: say the Mandarin for Beijing opera, then the Mandarin for a fan (of a singer).
\end{quote}

\subsection*{You'll want to know: xiōngdì}
\emph{xiōngdì} — \textquotedblleft{}brothers\textquotedblright{}.

\begin{quote}
\emph{xiōngdì liǎng ge gēge bǐ dìdi gāo} — Of the two brothers, the elder is taller.
\end{quote}

Its characters come later; for now it is a word for the ear.

\begin{cousinweb}
\emph{xiōngdì} is built from \emph{xiōng} 'an elder brother' and \emph{dì} 'a younger brother', both inherited from Middle Chinese, \emph{xjwaeng} and \emph{dejX} (Guangyun, in Baxter's transcription), which glosses the first, after an older dictionary, 'a man born first' and the second 'a man born after'.
\end{cousinweb}

\begin{grammarlens}[title={What you've built}]
One more word for saying how it was.
\end{grammarlens}

\subsection*{Your turn}
\begin{itemize}
\raggedright
  \item \emph{Say it:} \emph{xiōngdì}
  \item \emph{Say it:} \emph{xiōngdì}, once more
  \item \emph{Say it:} \emph{xiōngdì liǎng ge gēge bǐ dìdi gāo} — Of the two brothers, the elder is taller.
  \item \emph{Recall:} say the Mandarin for Beijing opera, then the Mandarin for a fan (of a singer), then say \emph{xiōngdì} again
\end{itemize}

\begin{tcolorbox}[breakable,colback=teal!4,colframe=teal!35!black,title={Before you move on}]
What does xiōngdì mean? (\textquotedblleft{}Brothers\textquotedblright{}.) Say it once more.
\end{tcolorbox}
\section[\texorpdfstring{\zh{父子} (fùzǐ)}{父子 (fùzǐ)}]{\zh{父子} (fùzǐ) — father and son}

\label{lesson:ZH-C507-fuzi}



\begin{quote}
Before the new one: say the Mandarin for a fan (of a singer), then the Mandarin for brothers.
\end{quote}

\subsection*{You'll want to know: \zh{父子}}
\textbf{\zh{父子}} — \emph{fùzǐ} — \textquotedblleft{}father and son\textquotedblright{}.

\begin{quote}
\textbf{\zh{他们父子都是歌手}} \emph{tāmen fùzǐ dōu shì gēshǒu} — Father and son are both singers.
\end{quote}

The new character, in pieces: \zh{父} — four strokes: two short strokes on top, one falling left and a dot falling right, then two long strokes that cross, one falling left and one falling right. The next lesson puts it in your hand.

\begin{cousinweb}
\textbf{\zh{父子}} (\emph{fùzǐ}) is built from \zh{父} 'a father' and \zh{子} 'a son', both inherited from Middle Chinese, \emph{bjuX} and \emph{tsiX} (Guangyun, in Baxter's transcription), which glosses the first, after an older dictionary, as the head of a family who leads and teaches, and the second 'offspring'. The dictionary also lists \zh{父} as \emph{pjuX}, a title of respect for a man.
\end{cousinweb}

\begin{grammarlens}[title={What you've built}]
One more word for saying how it was.
\end{grammarlens}

\subsection*{Your turn}
\begin{itemize}
\raggedright
  \item \emph{Say it:} \emph{fùzǐ}
  \item \emph{Say it:} \emph{fùzǐ}, once more
  \item \emph{Say it:} \emph{tāmen fùzǐ dōu shì gēshǒu} — Father and son are both singers.
  \item \emph{Recall:} say the Mandarin for a fan (of a singer), then the Mandarin for brothers, then say \emph{fùzǐ} again
\end{itemize}

\begin{tcolorbox}[breakable,colback=teal!4,colframe=teal!35!black,title={Before you move on}]
What does \zh{父子} mean? (\textquotedblleft{}Father and son\textquotedblright{}.) Say it once more.
\end{tcolorbox}
\section[\texorpdfstring{\zh{父} (fù)}{父 (fù)}]{\zh{父} — two short strokes, then two long ones}

\label{lesson:ZH-W507-fu-father}



\begin{quote}
The word first, then the character inside it.

\begin{itemize}
\raggedright
  \item \emph{Recall:} say \emph{fùzǐ} — \textbf{R1}, one lesson back, before it decays
\end{itemize}

\begin{quote}
\zh{人} \zh{子}
\end{quote}

\emph{Read it:} each one aloud, with nothing said first — \emph{rén}, \emph{zǐ}
\end{quote}

\subsection*{Script — two short strokes, then two long ones}
\hlblockfigure{\detokenize{figures/ZH-W507-fu-father-filmstrip.pdf}}{How \zh{父} is written, stroke by stroke}

\begin{quote}
\zh{父}
\end{quote}

\textbf{Two short strokes on top}: on the left a stroke falling to the left, on the right a dot falling to the right.

\textbf{Then two long strokes that cross}: a long stroke falling to the left, then a long stroke falling to the right, like the legs of \zh{人} set wide.

2 + 2 = four.

\subsection*{Your turn}
\begin{itemize}
\raggedright
  \item \emph{Write it:} \zh{父}, all four, with the model in front of you
  \item \emph{Cover:} write it again from memory, and count the pieces, not the strokes
\end{itemize}

\begin{tcolorbox}[breakable,colback=teal!4,colframe=teal!35!black,title={Before you move on}]
How many strokes does \zh{父} have? (\textbf{Four}.) Which way does the first stroke fall? (\textbf{To the left}.) What do the last two strokes do? (\textbf{They cross}.)

Source: \href{https://raw.githubusercontent.com/chanind/hanzi-writer-data/68d10a4b21150cae5e1ebbd223eed289cf32d90c/data/\%E7\%88\%B6.json}{Hanzi Writer Data \zh{父}, pinned stroke paths and medians}.
\end{tcolorbox}
\section[\texorpdfstring{\zh{父女} (fùnǚ)}{父女 (fùnǚ)}]{\zh{父女} (fùnǚ) — father and daughter}

\label{lesson:ZH-C507-funu}



\begin{quote}
Before the new one: say the Mandarin for brothers, then the Mandarin for a father and son.
\end{quote}

\subsection*{You'll want to know: \zh{父女}}
\textbf{\zh{父女}} — \emph{fùnǚ} — \textquotedblleft{}father and daughter\textquotedblright{}.

\begin{quote}
\textbf{\zh{他们父女一起去了舞会}} \emph{tāmen fùnǚ yìqǐ qù le wǔhuì} — Father and daughter went to the dance together.
\end{quote}

\begin{cousinweb}
\textbf{\zh{父女}} (\emph{fùnǚ}) is built from \zh{父} 'a father' and \zh{女} 'a daughter', both inherited from Middle Chinese, \emph{bjuX} and \emph{nrjoX} (Guangyun, in Baxter's transcription).
\end{cousinweb}

\begin{grammarlens}[title={What you've built}]
One more word for saying how it was.
\end{grammarlens}

\subsection*{Your turn}
\begin{itemize}
\raggedright
  \item \emph{Say it:} \emph{fùnǚ}
  \item \emph{Say it:} \emph{fùnǚ}, once more
  \item \emph{Say it:} \emph{tāmen fùnǚ yìqǐ qù le wǔhuì} — Father and daughter went to the dance together.
  \item \emph{Recall:} say the Mandarin for brothers, then the Mandarin for a father and son, then say \emph{fùnǚ} again
\end{itemize}

\begin{tcolorbox}[breakable,colback=teal!4,colframe=teal!35!black,title={Before you move on}]
What does \zh{父女} mean? (\textquotedblleft{}Father and daughter\textquotedblright{}.) Say it once more.
\end{tcolorbox}
\section[fūqī]{fūqī — married couple; husband and wife}

\label{lesson:ZH-C507-fuqi}



\begin{quote}
Before the new one: say the Mandarin for a father and son, then the Mandarin for a father and daughter.
\end{quote}

\subsection*{You'll want to know: fūqī}
\emph{fūqī} — \textquotedblleft{}married couple; husband and wife\textquotedblright{}.

\begin{quote}
\emph{tāmen zuò le sìshí nián fūqī} — They've been married for forty years.
\end{quote}

Its characters come later; for now it is a word for the ear.

\begin{cousinweb}
\emph{fūqī} is built from \emph{fū} 'a grown man' and \emph{qī} 'a wife', both inherited from Middle Chinese, \emph{pju} and \emph{tshej} (Guangyun, in Baxter's transcription), which glosses the first 'a grown man' and the second with a word of like sound meaning 'level'.
\end{cousinweb}

\begin{grammarlens}[title={What you've built}]
One more word for saying how it was.
\end{grammarlens}

\subsection*{Your turn}
\begin{itemize}
\raggedright
  \item \emph{Say it:} \emph{fūqī}
  \item \emph{Say it:} \emph{fūqī}, once more
  \item \emph{Say it:} \emph{tāmen zuò le sìshí nián fūqī} — They've been married for forty years.
  \item \emph{Recall:} say the Mandarin for a father and son, then the Mandarin for a father and daughter, then say \emph{fūqī} again
\end{itemize}

\begin{tcolorbox}[breakable,colback=teal!4,colframe=teal!35!black,title={Before you move on}]
What does fūqī mean? (\textquotedblleft{}Married couple; husband and wife\textquotedblright{}.) Say it once more.
\end{tcolorbox}
\section[zhǎngbèi]{zhǎngbèi — one's elders; the older generation}

\label{lesson:ZH-C507-zhangbei}



\begin{quote}
Before the new one: say the Mandarin for a father and daughter, then the Mandarin for a married couple.
\end{quote}

\subsection*{You'll want to know: zhǎngbèi}
\emph{zhǎngbèi} — \textquotedblleft{}one's elders; the older generation\textquotedblright{}.

\begin{quote}
\emph{zài zhǎngbèi miànqián wǒ gǎnjué hěn jǐnzhāng} — I feel nervous in front of my elders.
\end{quote}

Its characters come later; for now it is a word for the ear.

\begin{cousinweb}
\emph{zhǎngbèi} is built from \emph{zhǎng} 'elder' and \emph{bèi} 'a generation', both inherited from Middle Chinese, \emph{trjangX} and \emph{pwojH} (Guangyun, in Baxter's transcription), which glosses the first 'great' and the second 'people of one rank'.
\end{cousinweb}

\begin{grammarlens}[title={What you've built}]
One more word for saying how it was.
\end{grammarlens}

\subsection*{Your turn}
\begin{itemize}
\raggedright
  \item \emph{Say it:} \emph{zhǎngbèi}
  \item \emph{Say it:} \emph{zhǎngbèi}, once more
  \item \emph{Say it:} \emph{zài zhǎngbèi miànqián wǒ gǎnjué hěn jǐnzhāng} — I feel nervous in front of my elders.
  \item \emph{Recall:} say the Mandarin for a father and daughter, then the Mandarin for a married couple, then say \emph{zhǎngbèi} again
\end{itemize}

\begin{tcolorbox}[breakable,colback=teal!4,colframe=teal!35!black,title={Before you move on}]
What does zhǎngbèi mean? (\textquotedblleft{}One's elders; the older generation\textquotedblright{}.) Say it once more.
\end{tcolorbox}
